\chapter{The Huffman Dictionaries}
\label{app:huffman}

Table~\ref{tab:huffman} gives the complete original dictionary
\cite[Table~A.1]{kiem2025}, used in Stage~1, in the no-prediction baseline and
in the original-dictionary board builds. The published codes are most-significant-bit first; the
implementation packs least-significant-bit first, so the stored codewords are
bit-reversed. Both forms are listed, and all sixteen entries were verified to
agree with the published table.

The reversed code set remains prefix-free, which is what allows the decoder to
match candidates from the low end of its bit buffer with no ambiguity.

\begin{table}[h]
\centering
\caption[The original Huffman dictionary]{The original Huffman dictionary. ``Cost'' is
the total bits to code one residual in that category, $\ell + S_4$, from
Equation~\eqref{eq:cost}.}
\label{tab:huffman}
\small
\begin{tabular}{@{}rlrlrr@{}}
\toprule
$S_4$ & $|v|$ range & $\ell$ & code (MSB-first) & stored (LSB-first) & cost \\
\midrule
0  & $v = 0$          & 4  & \texttt{1010}           & \texttt{0x0005} & 4 \\
1  & 1                & 3  & \texttt{100}            & \texttt{0x0001} & 4 \\
2  & 2--3             & 2  & \texttt{01}             & \texttt{0x0002} & 4 \\
3  & 4--7             & 2  & \texttt{00}             & \texttt{0x0000} & 5 \\
4  & 8--15            & 2  & \texttt{11}             & \texttt{0x0003} & 6 \\
5  & 16--31           & 5  & \texttt{10110}          & \texttt{0x000D} & 10 \\
6  & 32--63           & 6  & \texttt{101110}         & \texttt{0x001D} & 12 \\
7  & 64--127          & 7  & \texttt{1011110}        & \texttt{0x003D} & 14 \\
8  & 128--255         & 8  & \texttt{10111110}       & \texttt{0x007D} & 16 \\
9  & 256--511         & 9  & \texttt{101111110}      & \texttt{0x00FD} & 18 \\
10 & 512--1023        & 10 & \texttt{1011111110}     & \texttt{0x01FD} & 20 \\
11 & 1024--2047       & 11 & \texttt{10111111110}    & \texttt{0x03FD} & 22 \\
12 & 2048--4095       & 14 & \texttt{10111111111111} & \texttt{0x3FFD} & 26 \\
13 & 4096--8191       & 14 & \texttt{10111111111110} & \texttt{0x1FFD} & 27 \\
14 & 8192--16383      & 13 & \texttt{1011111111110}  & \texttt{0x0FFD} & 27 \\
15 & 16384--32767     & 12 & \texttt{101111111110}   & \texttt{0x07FD} & 27 \\
\bottomrule
\end{tabular}
\end{table}

\noindent
Note the non-monotonicity at categories 12--15: the longest codewords are
assigned to categories 12 and 13 rather than to 15. This is a property of the
published dictionary and is reproduced faithfully. It has no effect on this
dataset, where no residual exceeds category 10.


\section*{The re-derived dictionaries}

Table~\ref{tab:rederived} gives the codeword lengths of the dictionaries
re-derived for each algorithm (Section~\ref{sec:dict}), built from the $S_4$
counts of frames 0--4 with one added to every category, and the total cost of
each category with each. The codewords themselves are the canonical Huffman
codes for these lengths, bit-reversed for the packer; they are the tables
\texttt{DRHE\_TDM\_HUFFMAN\_TABLE} and \texttt{LPC\_TDM\_HUFFMAN\_TABLE} in the
HLS source.

\begin{table}[h]
\centering
\caption[The re-derived dictionaries]{Codeword length $\ell$ and total cost
$\ell + S_4$ per category for the original and the two re-derived dictionaries.}
\label{tab:rederived}
\small
\begin{tabular}{@{}r rr rr rr@{}}
\toprule
& \multicolumn{2}{c}{original} & \multicolumn{2}{c}{DRHE} & \multicolumn{2}{c}{LPC}\\
\cmidrule(lr){2-3}\cmidrule(lr){4-5}\cmidrule(lr){6-7}
$S_4$ & $\ell$ & cost & $\ell$ & cost & $\ell$ & cost\\
\midrule
0 & 4 & 4 & 3 & 3 & 3 & 3\\
1 & 3 & 4 & 2 & 3 & 1 & 2\\
2 & 2 & 4 & 1 & 3 & 2 & 4\\
3 & 2 & 5 & 4 & 7 & 4 & 7\\
4 & 2 & 6 & 5 & 9 & 5 & 9\\
5 & 5 & 10 & 6 & 11 & 6 & 11\\
6 & 6 & 12 & 7 & 13 & 7 & 13\\
7 & 7 & 14 & 8 & 15 & 8 & 15\\
8 & 8 & 16 & 9 & 17 & 9 & 17\\
9 & 9 & 18 & 10 & 19 & 10 & 19\\
10 & 10 & 20 & 11 & 21 & 13 & 23\\
11 & 11 & 22 & 14 & 25 & 13 & 24\\
12 & 14 & 26 & 14 & 26 & 13 & 25\\
13 & 14 & 27 & 13 & 26 & 13 & 26\\
14 & 13 & 27 & 13 & 27 & 12 & 26\\
15 & 12 & 27 & 13 & 28 & 12 & 27\\
\bottomrule
\end{tabular}
\end{table}

\noindent
The re-derived dictionaries code their commonest category in three bits (DRHE,
category 2) or two (LPC, category 1), where the original cannot code anything
in fewer than four. They are longer above category 3, which costs little
because those categories are rare after prediction.

\chapter{Bitstream Layouts}
\label{app:bitstream}

\section{DRHE}

DRHE's bitstream is a single continuous sequence of coded residuals, in the
raster order the samples arrive:
\[
\underbrace{\text{enc}(d_{\Re})\ \text{enc}(d_{\Im})}_{\text{ramp }0,\ \text{bin }0,\ \text{ch }0}\
\underbrace{\text{enc}(d_{\Re})\ \text{enc}(d_{\Im})}_{\text{ramp }0,\ \text{bin }0,\ \text{ch }1}\
\cdots
\]
where $\text{enc}(v)$ is the Huffman codeword for $S_4(v)$ followed by
$S_4(v)$ APPEND bits. There is no header: the frame dimensions arrive over the
AXI4-Lite control interface. The final AXI word is zero-padded.

\section{LPC}

LPC's bitstream has two sections:
\[
\underbrace{\big[a_1^{\Re}\ a_2^{\Re}\ a_1^{\Im}\ a_2^{\Im}\big]_{\text{ch},\text{bin}}}_{\text{coefficient table, }4\times16\text{ bits each}}
\quad
\underbrace{\text{enc}(d_{\Re})\ \text{enc}(d_{\Im})\ \cdots}_{\text{residuals, raster order}}
\]
The coefficient table comes first because the decoder needs all of it before it
can reconstruct anything. The residuals follow in the original raster order ---
not in the order the compressor computed them --- which is what allows the
decompressor to operate without a frame buffer
(Section~\ref{sec:lpcimpl}).


\chapter{Reproducing the Results}
\label{app:repro}

All sources are in the thesis repository under \texttt{ThesisA/}. The build
tree must sit on a path without spaces, which Vitis rejects.

\section{Source files}

\begin{longtable}{@{}p{58mm}p{82mm}@{}}
\toprule
File & Purpose \\
\midrule
\multicolumn{2}{@{}l}{\textbf{MATLAB} (\texttt{Matlab\_Sim/})}\\
\texttt{loadColoRadarFrame.m} & reads a frame and applies the TDM flattening \\
\texttt{preprocessing.m}, \texttt{rangeFFT.m} & DC removal, window, range FFT \\
\texttt{compress\_fx16.m} & FX16 quantisation, Equation~\eqref{eq:scale} \\
\texttt{compress\_drhe.m}, \texttt{decompress\_drhe.m} & DRHE; lag from \texttt{config.predLag}, default $\nTx$ \\
\texttt{compress\_lpc\_huffman.m}, \texttt{decompress\_lpc\_huffman.m} & LPC at lags $L$, $2L$ \\
\texttt{huffman\_tables.m}, \texttt{huffman\_code\_lengths.m} & the original and re-derived dictionaries; retraining \\
\texttt{run\_coloradar\_benchmark.m} & Stage~1: all eight methods with detection metrics \\
\texttt{make\_poster\_figs.m} & Stage~2 breakdown, dictionary tables, and the figures of Chapter~\ref{ch:results}; parameters at the top \\
\texttt{export\_frames\_for\_vitis.m} & writes the \texttt{int16} stimulus for the HLS testbenches \\
\midrule
\multicolumn{2}{@{}l}{\textbf{HLS} (\texttt{hls\_component/})}\\
\texttt{drhe\_tdm\_common.h}, \texttt{drhe\_tdm\_compress.cpp}, \texttt{drhe\_tdm\_decompress.cpp} & DRHE cores and the DRHE dictionary \\
\texttt{lpc\_tdm\_common.h}, \texttt{lpc\_tdm\_compress.cpp}, \texttt{lpc\_tdm\_decompress.cpp} & LPC cores, Yule--Walker solver and the LPC dictionary \\
\texttt{*\_tb.cpp}, \texttt{*\_tb\_edge.cpp} & 50-frame and edge-case testbenches \\
\midrule
\multicolumn{2}{@{}l}{\textbf{Board} (\texttt{fpga\_kcu116/})}\\
\texttt{hls/export\_ips.bat} & C-simulation (writes the golden streams) and IP packaging \\
\texttt{vivado/build.bat} & block design, synthesis, implementation, bitstream; \texttt{decomp 1} adds the decompressor \\
\texttt{sw/src/main.c}, \texttt{board\_map.h} & MicroBlaze firmware \\
\texttt{host/run\_board.bat}, \texttt{run\_board.tcl} & program the board, load frames, run, read results \\
\texttt{host/pc\_decompress.bat} & PC decompression with the HLS decoder source \\
\texttt{host/compare\_results.py} & byte comparison with C-simulation, reconstruction check, tables \\
\midrule
\multicolumn{2}{@{}l}{\textbf{Independent check} (\texttt{verify\_indep/}, Python)}\\
\texttt{indep.py} & the whole chain from raw frame to bit count \\
\texttt{retrain\_complete.py} & dictionaries trained on frames 0--4, tested on 5--49 \\
\texttt{golden\_check.py}, \texttt{independent\_decode.py} & residual-level check and decoding of the hardware bitstreams \\
\bottomrule
\end{longtable}

\section{Commands}

C-simulation, synthesis and co-simulation of a core (DRHE shown):
\begin{verbatim}
vitis-run --mode hls --csim  --config test_tdm.cfg --work_dir hls_tdm
v++ -c    --mode hls         --config test_tdm.cfg --work_dir hls_tdm
vitis-run --mode hls --cosim --config test_tdm.cfg --work_dir hls_tdm
\end{verbatim}

\noindent
Board run, from \texttt{fpga\_kcu116/host} (\texttt{<tag>} is \texttt{drhe},
\texttt{lpc}, \texttt{drhe\_full} or \texttt{lpc\_full}):
\begin{verbatim}
run_board.bat <tag> 50
pc_decompress.bat <tag>          (compress-only builds)
python compare_results.py <tag>
\end{verbatim}

\noindent
Figures and tables of Chapter~\ref{ch:results}, in MATLAB from
\texttt{Matlab\_Sim}: \texttt{make\_poster\_figs}. Simulation results are cached;
set \texttt{P.recompute = true} after changing the algorithm code.

\chapter{Verification Record}
\label{app:verify}

Table~\ref{tab:verify} lists every check on which the results rest, where it
was made, and its outcome.

\begin{longtable}{@{}p{55mm}p{38mm}p{45mm}@{}}
\caption[Verification record]{The verification record.}\label{tab:verify}\\
\toprule
Check & Where & Outcome\\
\midrule
\endfirsthead
\toprule
Check & Where & Outcome\\
\midrule
\endhead
Stage~1, eight methods, 50 frames & MATLAB & Table~\ref{tab:stage1}; re-run reproduces every cell\\
DRHE and LPC lossless, original dictionary & HLS C-sim, 50 frames & MaxDiff 0 on every frame\\
DRHE and LPC lossless, re-derived dictionaries & HLS C-sim, 50 frames & MaxDiff 0; CR \num{4.21704}, \num{4.40345}\\
MATLAB against HLS bit counts & both & equal up to the last 256-bit padding word\\
Independent bit counts, original dictionary & Python, 50 frames & equal to HLS on every frame\\
Independent decode, re-derived dictionaries & Python & no residual mismatch\\
Edge cases, eight synthetic inputs & HLS C-sim & all pass within their allowed MaxDiff\\
RTL behaviour & co-simulation, 1 frame & MaxDiff 0, both algorithms\\
Timing closure, four builds & Vivado & all met at \SI{100}{\mega\hertz} (Table~\ref{tab:builds})\\
Compressed bytes against C-sim, original dictionary & KCU116, 50 frames & byte-identical, both algorithms\\
Decompression of board output & PC, 50 frames & MaxDiff 0, both algorithms\\
Compressed bytes against C-sim, re-derived dictionaries & KCU116, 50 frames & byte-identical; CR \num{4.21704}, \num{4.40345}\\
On-chip decompression (\texttt{\_full}) & KCU116, 50 frames & MaxDiff 0, both algorithms; \SI{1.97}{\milli\second} and \lpcDecTime{} per frame\\
Parallel LPC decoder against the first one & C-sim, 50 frames + 8 edge cases & bit-identical output\\
Frames ending on a word boundary & C-sim, 8 edge cases, both algorithms & TLAST always sent; ColoRadar output unchanged\\
DRHE matches Kiem's specification & rebuilt simulation & bit-identical output\\
\bottomrule
\end{longtable}
